\documentclass[11pt]{article}
\usepackage[margin=1in]{geometry}
\usepackage{amsmath,amssymb,amsthm}
\usepackage[hidelinks]{hyperref}
\newtheorem{lemma}{Lemma}
\newtheorem{theorem}[lemma]{Theorem}
\newtheorem{corollary}[lemma]{Corollary}
\newtheorem{proposition}[lemma]{Proposition}
\newcommand{\Cp}{\mathbb C_p}
\newcommand{\Zp}{\mathbb Z_p}
\newcommand{\Q}{\mathbb Q}
\newcommand{\C}{\mathbb C}
\newcommand{\OO}{\mathcal O}
\title{Algebraic Gamma shapes in the prime-power arithmetic domains}
\author{}\date{}
\begin{document}
\maketitle

Fix a prime $p$ and a positive integer $s$, and put
\[
 q=p^s,\qquad b_s=\frac{p^s-1}{p-1}.
\]
Let $\Cp$ be a complete algebraically closed extension of $\mathbb Q_p$,
with valuation $v(p)=1$ and valuation ring $\OO=\{x:v(x)\ge0\}$.
The valuation need not be integer-valued.
Define the translation-stable domains
\[
 \Omega_{p,1}=\OO,\qquad
 \Omega_{p,s}=\Zp+p^s\OO\quad(s\ge2).
 \tag{1}
\]
For $s\ge2$ the second set is the finite union of closed disks
\[
 \Omega_{p,s}=\bigcup_{m=0}^{q-1}(-m+p^s\OO).
 \tag{2}
\]
We use
\[
 T_{n,\gamma,q}(t)=\sum_{k=0}^n\binom nk(\gamma)_{qk}t^k
\]
and the integer-coefficient cyclic determinant
$\mathcal D_{n,q}(\gamma)$ of the
\href{ARITHMETIC_FRONTIER_TERMINAL_ORDERS.tex}{retained arithmetic proof}.
The point of this note is to record which algebraic parameters that proof
covers. It does not extend its proper-prime-power assertion to all
of $\OO$.

\begin{lemma}[The local determinant domains]\label{lem:local-domains}
For every integer $n\ge1$, put $N=\binom n2$.
At prime gaps the coefficientwise congruence is
\[
 \mathcal D_{n,p}(X)\equiv(-1)^N\pmod p
       \quad\text{in }\mathbb F_p[X].
 \tag{3}
\]
Consequently, for every $\gamma\in\OO$, with any residue field,
$v(\mathcal D_{n,p}(\gamma))=0$.
For every $s\ge2$ and every $\gamma\in\Omega_{p,s}$,
\[
 v(\mathcal D_{n,q}(\gamma))=(b_s-s)N.
 \tag{4}
\]
In particular the determinant is nonzero throughout the stated domains.
\end{lemma}
\begin{proof}
The coefficientwise assertion (3) is exactly
\texttt{thm:prime-gap-unit-terminal} of the arithmetic proof.
It is stronger than a statement only about values at $\Zp$.
Indeed it writes the integer polynomial as
$(-1)^N+pH(X)$ with $H\in\mathbb Z[X]$.
For $\gamma\in\OO$, $H(\gamma)\in\OO$, so evaluation is a unit.
The proof of (3) retains the polynomial
$u+\gamma^p-\gamma$; it does not substitute $\gamma^p=\gamma$.

For $s\ge2$, the proof of
\texttt{thm:prime-power-lower-one} first excludes determinant roots
on every disk in (2), working in a complete algebraically closed valued
extension. The restriction to $\gamma\in\Zp$ occurs only at its final
covering step. We check this distinction explicitly.

If $\gamma=-m+p^s\eta$, $\eta\in\OO$, then at each integer level
$1\le r\le s$ exactly $qk/p^r$ of the factors
$\gamma,\ldots,\gamma+qk-1$ have valuation at least $r$.
The remaining valuations are nonnegative, so
\[
 v((\gamma)_{qk})\ge \sum_{r=1}^s qk/p^r=b_sk.
 \tag{5}
\]
This lower bound remains valid for fractional valuations.
The monic reciprocal
\[
 R(z)=z^{qn}T_{n,\gamma,q}((-1)^{q+1}z^{-q})
\]
therefore has every nonzero root in $v(z)\ge b_s/q$:
at a point of smaller valuation its leading term is uniquely least.
A zero of the determinant, by the division-free cyclic presentation,
the no-origin lemma and the critical-ideal identity used in the arithmetic
proof, would give a nonzero cyclic argument $u=z^q$ with $v(u)\ge b_s$.

The actual operator has constant
$c=\gamma+q-1\in\Zp+p^s\OO$ and seed $1$.
It lies in the exact class of the arithmetic proof's
\texttt{lem:artin-hasse-closure} and
\texttt{lem:frobenius-noncollision}.
Those lemmas allow coefficients in $\OO$ and use its maximal ideal;
they do not require a discrete valuation.
Their series and every finite Euler derivative converge at $v(z)>0$,
since their coefficients are integral. The radius bound above is strictly
positive. The primitive-polynomial normalization in the descent uses an
actual coefficient of minimum valuation, so it does not assume that
the minimum is an integer or a power of $p$ available in the base field.
The resulting contradiction excludes determinant roots on each whole
closed disk, including its algebraic extension points.

For completeness, the center calculation in the same proof is
\[
 \mathcal D_{n,q}(-m)=
 \frac{((qn-m-1)!)^{n-1}}
 {q^N\prod_{k=1}^{n-1}k^k(qk-m-1)!},
 \qquad
 v(\mathcal D_{n,q}(-m))=(b_s-s)N.
 \tag{6}
\]
The case $n=1$ uses empty products.
Factor $\mathcal D$ over $\Cp$. Every root $\rho$ lies outside the
closed disk, so $v(\rho+m)<s$.
For every $\gamma$ in that disk,
$v(\gamma-\rho)=v(-m-\rho)$ by the strict ultrametric inequality.
Thus the polynomial has constant valuation on the entire disk,
equal to (6). This proves (4), not just its restriction to $\Zp$.
\end{proof}

\begin{theorem}[Positive algebraic shapes]\label{thm:algebraic-lower-one}
Let $\beta\in\overline{\Q}\cap\mathbb R_{>0}$ admit an embedding
\[
 \iota:\Q(\beta)\longrightarrow\Cp,\qquad
                  \iota(\beta)\in\Omega_{p,s}.
 \tag{7}
\]
For every $(A,B)\in\C^2\setminus\{(0,0)\}$ and
$f(X)=X(A+BX^q)$, the Gamma functional
$\mathcal L_\beta(X^j)=(\beta)_j$ satisfies
\[
 \bigl(\mathcal L_\beta(f^m),\mathcal L_\beta(f^{m+1})\bigr)
                     \ne(0,0)\qquad(m\ge1).
 \tag{8}
\]
For prime gaps, it suffices that $\beta$ be integral at one $p$-adic
embedding; its residue need not lie in $\mathbb F_p$.
For every prime-power gap, a $p$-integral embedding into $\mathbb Q_p$
suffices. The disk condition (7) also permits embeddings into larger,
including ramified, extensions.
\end{theorem}
\begin{proof}
For terminal order $n=m+1\ge2$, set $\gamma=m+\beta>0$.
Integer translation preserves (1), so $\iota(\gamma)$ belongs to the
same domain. Lemma~\ref{lem:local-domains} gives
$\mathcal D_{n,q}(\gamma)\ne0$ in $\Q(\beta)$:
a field embedding cannot send a nonzero value to zero, or a zero value
to a nonzero value.
The all-gap discriminant identity used in the arithmetic proof is
\[
 \operatorname{Disc} T_{n,\gamma,q}
 =(-1)^Nq^N\left(\prod_{k=1}^n k^k\right)
       \left(\prod_{k=1}^{n-1}(\gamma)_{qk}\right)
       \mathcal D_{n,q}(\gamma).
 \tag{9}
\]
Every rising factorial on the right is nonzero because the chosen real
shape $\gamma$ is positive. Hence this algebraic discriminant is nonzero,
and $T_{n,\gamma,q}$ is squarefree after extension to $\C$.

If $AB\ne0$, the normalized two moments are
$T_{n,\gamma+1,q}(B/A)$ and $T_{n-1,\gamma,q}(B/A)$.
Writing $T=T_{n,\gamma,q}$ and $\theta=tD_t$, coefficient comparison gives
\[
 T_{n,\gamma+1,q}=T+(q/\gamma)\theta T,\qquad
 T_{n-1,\gamma,q}=T-(1/n)\theta T.
\]
The coefficient determinant is
$-(\gamma+qn)/(\gamma n)\ne0$.
A common zero would therefore be a multiple zero of $T$, impossible;
$t=0$ is excluded by $T(0)=1$.
The normalization factors $A^j(\beta)_j$ are nonzero.
If $A=0$ or $B=0$, the phase is a nonzero monomial and every positive
moment is directly nonzero.
\end{proof}


\begin{corollary}[Prime gaps at algebraic shapes]\label{cor:prime-gap-algebraic-shapes}
Let $p$ be prime and let $\beta>0$ be algebraic over $\Q$.
If $v_{\mathfrak p}(\beta)\ge0$ at at least one place
$\mathfrak p$ of $\Q(\beta)$ above $p$, then every nonzero
\[
 f(X)=X(A+BX^p),\qquad A,B\in\C,
\]
has a nonzero Gamma moment in every pair of consecutive positive orders.
In particular every positive algebraic-integer shape is covered at
every prime gap, including $\beta=\sqrt2$.
\end{corollary}
\begin{proof}
A place above $p$ is realized by an embedding of $\Q(\beta)$ into $\Cp$.
For $s=1$, its integral parameter lies in $\Omega_{p,1}=\OO$.
At terminal order $n$, the parameter $\gamma=n+\beta-1$
remains integral there. Thus (3) evaluates to a unit, and
Theorem~\ref{thm:algebraic-lower-one} applies.
The argument needs only one such embedding and allows arbitrary
complex $A,B$; a field embedding certifies nonvanishing of the algebraic
discriminant independently of those coefficients.

Every algebraic integer is integral at every such embedding.
Indeed, at a root of a monic integer polynomial, a negative valuation
would make its leading term uniquely least among all terms.
This is impossible. No unit condition on $\beta$ is needed:
at $p=2$ the valuation of $\sqrt2$ is $1/2$, which is allowed.
\end{proof}

\begin{corollary}[Necessary shape locus for a prime-gap exception]
\label{cor:prime-gap-exception-locus}
Fix a prime $p$ and a positive real shape $\beta$.
If some nonzero $X(A+BX^p)$ has two consecutive positive
Gamma moments zero, then $\beta$ is algebraic and
\[
 v_{\mathfrak p}(\beta)<0
 \quad\text{at every place }\mathfrak p\mid p
 \text{ of }\Q(\beta).
 \tag{11}
\]
For the primitive integer minimal polynomial
$g_\beta(X)=a_0+\cdots+a_dX^d$, condition (11) is equivalent to
\[
 \overline{g_\beta}\in\mathbb F_p^\times;
 \quad\text{that is, }p\nmid a_0
 \text{ and }p\mid a_j\text{ for every }j\ge1.
 \tag{12}
\]
These are necessary conditions for an exception, not sufficient
conditions for failure.
\end{corollary}
\begin{proof}
The \href{BINOMIAL_IRREDUCIBILITY_REDUCTION.tex}{transcendental-shape theorem},
Corollary \texttt{cor:transcendental-shape}, already gives all-index
consecutive coprimality for every positive transcendental shape and
every binomial support. A failure here therefore forces $\beta$ to
be algebraic. Corollary~\ref{cor:prime-gap-algebraic-shapes} then
forces every place above $p$ to be nonintegral. This proves (11).

For completeness we prove the elementary equivalence with (12).
If (12) holds, evaluation of $g_\beta$ at any integral point of $\Cp$
is a unit, so it has no integral root.
Conversely assume all its $p$-adic roots $\beta_i$ have negative
valuation; this is exactly (11), since all embeddings above a given
place give the same valuation sign.
Because $\beta\ne0$, the constant coefficient $a_0$ is nonzero.
The reciprocal roots all have positive valuation, and
\[
 g_\beta(X)=a_0\prod_{i=1}^d(1-X/\beta_i).
\]
Vieta's formulas give $v_p(a_j/a_0)>0$ for every nonzero
coefficient with $j\ge1$. These are ratios of integers, so
$v_p(a_j)\ge v_p(a_0)+1$.
Primitiveness forces the minimum coefficient valuation to be zero;
hence $v_p(a_0)=0$ and all other coefficients are divisible by $p$.
This proves (12).

The larger integral domain here is a consequence of the
\emph{prime-gap} coefficient congruence.
At proper prime powers, the sufficient domains in (1)--(2)
remain unchanged; no assertion about all algebraic-integer shapes
at every prime-power gap follows.
\end{proof}

\begin{corollary}[Fixed monomial radial markers]\label{cor:algebraic-marked-radial}
Let independent $Y_i\sim\operatorname{Gamma}(\alpha_i,1)$ have
arbitrary positive real shapes, $1\le i\le r$, $r\ge1$.
Suppose their total shape $b=\sum_i\alpha_i$ is positive algebraic
and has an embedding satisfying (7).
Set
\[
 S=\sum_iY_i,\qquad F=S(A+BS^q),\qquad (A,B)\ne(0,0).
\]
For every fixed $\nu\in\mathbb N_0^r$ and every $m\ge1$,
\[
 \bigl(\mathbb E[Y^\nu F^m],\mathbb E[Y^\nu F^{m+1}]\bigr)\ne(0,0).
\]
Individual shapes need not be algebraic.
\end{corollary}
\begin{proof}
The exact monomial tilt in
\href{PRIME_POWER_MARKED_RADIAL_SFC.tex}{the marked radial identity}
gives
\[
 \mathbb E[Y^\nu h(S)]
  =\left(\prod_i(\alpha_i)_{\nu_i}\right)
           \mathcal L_{b+|\nu|}(h).
\]
The shifted shape is positive algebraic, and its embedding is still
in (1). The prefactor is strictly positive. Apply
Theorem~\ref{thm:algebraic-lower-one} to the same fixed marker in
both entries.
\end{proof}

The disk condition also allows higher-degree local embeddings:
for example, the positive algebraic shape $1+p^s\sqrt2$ lies in
$\Zp+p^s\OO$ at every embedding, whether or not that embedding
takes values in $\mathbb Q_p$. These sufficient domains are not
claimed to classify all positive algebraic shapes.

\begin{proposition}[A residue-extension obstruction]\label{prop:residue-extension-obstruction}
At $q=4,n=2$, extending the local determinant or squarefreeness
conclusion to every integral algebraic parameter is false.
There is an algebraic $2$-adic unit $\gamma$ whose residue has degree two
over $\mathbb F_2$ and for which
\[
 \mathcal D_{2,4}(\gamma)=0,\qquad
                 T_{2,\gamma,4}(t)=(1-45t)^2.
\]
This does not give a positive-real-Gamma counterexample.
\end{proposition}
\begin{proof}
The degree-two discriminant identity gives, polynomially,
\[
 \mathcal D_{2,q}(X)=\frac{(X+q)_q-(X)_q}{q}.
\]
For $q=4$ this is
\[
 \mathcal D_{2,4}(X)=2(2X+7)(X^2+7X+15),\qquad
 \frac{\mathcal D_{2,4}(X)}2\equiv X^2+X+1\pmod2.
 \tag{10}
\]
Take either root $\gamma$ of $X^2+7X+15$ in $\Cp$.
The monic polynomial and unit constant term force $v(\gamma)=0$.
Its reduction satisfies the irreducible polynomial
$X^2+X+1$ over $\mathbb F_2$, so it is outside the prime-field residue
class and outside every disk in (2).
Reduction modulo $X^2+7X+15$ gives
$(\gamma)_4=-45$ and $(\gamma)_8=2025$, proving the displayed square.
The complex roots of this quadratic are
$(-7\pm\sqrt{-11})/2$, both nonreal. Thus this obstruction concerns
the larger arithmetic parameter domain, not the positive-shape
assertion (8).
\end{proof}

The prime-gap coefficient congruence is uniform on all integral
algebraic parameters. At proper prime powers, the proved integer
closed disks remain the sufficient domain; neither an assertion about
all integral extension points nor an assertion about every positive
algebraic shape is inferred. No novelty, external review or formal
verification is claimed.
\end{document}
